
\documentclass[11pt]{article}

\usepackage[margin=1in]{geometry}
\usepackage{amsmath,amssymb,amsfonts}
\usepackage{bm}
\usepackage{mathtools}
\usepackage{booktabs}
\usepackage{enumitem}
\usepackage{graphicx}
\usepackage{subcaption}
\usepackage{float}
\usepackage{algorithm}
\usepackage{algpseudocode}
\usepackage[hidelinks]{hyperref}
\usepackage{multirow}
\hypersetup{breaklinks=true}
\emergencystretch=2em
\numberwithin{equation}{section}
\captionsetup[subfigure]{labelformat=simple}
\renewcommand\thesubfigure{\Alph{subfigure}}

\graphicspath{{figs/}}

\title{AHBA--GTEx Imputation by Deterministic and Probabilistic Latent Alignment}
\author{}
\date{}

\begin{document}
\maketitle

\section{Introduction}
The Allen Human Brain Atlas (AHBA) and GTEx provide complementary but structurally
mismatched views of human brain transcriptomics.  AHBA supplies dense atlas-level
measurements over a canonical parcellation, whereas GTEx provides subject-specific
observations over only a sparse subset of regions \cite{hawrylycz2012anatomically}.
GTEx broadens population coverage and subject diversity, but only over a limited set
of brain regions \cite{gtex2013genotype,gtex2017genetic,gtex2020atlas}.  Our aim is
to impute, for each GTEx subject, a complete parcel-by-gene expression matrix over
all canonical parcels by leveraging the shared spatial structure encoded in the AHBA
atlas.  This is not a standard missing-data problem: it is cross-platform
missing-region inference under domain shift, because the atlas reference is derived
from microarray measurements (AHBA) while subject-level observations in GTEx arise
from bulk RNA sequencing.

A common assumption in this setting is that inter-regional transcriptomic variation
is dominated by a low-dimensional spatial structure that is more completely visible
in the atlas than in any individual GTEx subject.  The central modeling question is
therefore how to infer that structure, align sparse subject measurements to the atlas
frame, and extend subject-specific expression over unobserved parcels.  Two
complementary approaches are developed here.  The first is the deterministic latent
alignment model (DLAM), based on subject-specific partial least squares (PLS), affine
latent alignment, and spatial interpolation.  The second is the probabilistic latent
alignment model (PLAM), in which the atlas defines a global factor structure, each
subject is represented by an orthogonally aligned latent deviation field, and
unobserved parcels are inferred by Gaussian-process conditioning.

\section{Data Setting and Notation}
Let $R=150$ denote the canonical parcel set and $G$ the number of genes under
analysis.  AHBA defines a dense parcel-by-gene atlas
\[
  X_A \in \mathbb{R}^{R \times G},
\]
while subject $s$ contributes GTEx observations on a sparse subset of parcels
\[
  \Omega_s \subset \{1,\dots,R\},
  \qquad
  X_s^{(G)} \in \mathbb{R}^{|\Omega_s| \times G}.
\]
All spatial quantities are defined with respect to canonical parcel centroids
\[
  Y \in \mathbb{R}^{R \times 3}.
\]
The latent dimension is denoted by $K$ and fixed to $K=3$ throughout.  Atlas-scale
heatmaps and scatter diagnostics are shown on the matched high-variance subset with
$G=88$ for readability, whereas leave-one-region-out diagnostics are computed over the
full gene set so that the cross-validated evaluation remains aligned with
transcriptome-wide completion.

\section{Cross-Dataset Harmonization}
\label{sec:harmonization}

Both alignment methods begin with the same preprocessing step: AHBA and GTEx are
mapped into a common harmonized expression domain before any latent alignment is
attempted.  The purpose of harmonization is to remove systematic, platform-induced
differences in marginal gene-expression distributions while preserving the relative
biological variation across parcels and subjects.

\subsection{General harmonization model}
\label{sec:harmonization-general}

Let $\mathcal{H}_A$ and $\mathcal{H}_G$ be platform-specific transformation functions
parameterized by shared parameters $\theta_H$, satisfying the following two
properties: (i)~the mapped expressions are on a common scale, in the sense that
$\mathcal{H}_A(X_A;\theta_H)$ and $\mathcal{H}_G(X_s^{(G)};\theta_H)$ have compatible
marginal distributions across genes; and (ii)~biologically informative between-parcel
and between-subject contrasts are preserved under the transformations.  The shared
preprocessing transformation is written as
\begin{equation}
  \label{eq:shared-harmonization}
  X_A^{(h)} = \mathcal{H}_A(X_A;\theta_H),
  \qquad
  X_{s,h}^{obs} = \mathcal{H}_G(X_s^{(G)};\theta_H).
\end{equation}
The parameters $\theta_H$ are estimated using the AHBA atlas together with the
available GTEx rows, treating the two platforms as two batches.  Both DLAM and PLAM
receive the harmonized atlas
$X_A^{(h)}$ and harmonized subject observations $X_{s,h}^{obs}$ as input, and both
ultimately produce a completed harmonized subject atlas
\[
  \hat{X}_s^{(h)} \in \mathbb{R}^{R\times G}.
\]

\subsection{Implementation via ComBat-style Harmonization}
\label{sec:harmonization-combat}

In the default pipeline, Eq.~\eqref{eq:shared-harmonization} is implemented by a
ComBat-style batch harmonizer fit jointly to AHBA and eligible GTEx rows
\cite{johnson2007combat}.  We follow the notation of Johnson et~al. and write
$y_{ijg}$ for the observed expression of gene $g$ in sample $j$ from batch $i$.
In our setting, the two batches are $i=A$ for AHBA and $i=G$ for GTEx.  ComBat
models each gene as
\begin{equation}
  \label{eq:combat-model}
  y_{ijg}
  =
  \alpha_g
  +
  \mathbf{x}_{ij}^{\top}\boldsymbol{\beta}_g
  +
  \gamma_{ig}
  +
  \delta_{ig}\varepsilon_{ijg},
\end{equation}
where $\alpha_g$ is the gene-specific intercept, $\mathbf{x}_{ij}$ is a vector of
row-level covariates, $\boldsymbol{\beta}_g$ gives the gene-specific covariate
effects, $\gamma_{ig}$ is the additive batch effect, and $\delta_{ig}$ is the
multiplicative batch scale.  In the default analysis, $\mathbf{x}_{ij}$ contains
normalized age, a sex indicator, and broad anatomical macro-system indicators, with
cortical association cortex used as the reference macro-system.

Operationally, the pooled AHBA--GTEx expression table is first adjusted for the
fitted covariate contribution $\mathbf{x}_{ij}^{\top}\hat{\boldsymbol{\beta}}_g$.
For each gene, the covariate-adjusted values are centered and scaled using pooled
statistics across both platforms:
\[
  s_{ijg}
  =
  \frac{
    y_{ijg}
    -
    \mathbf{x}_{ij}^{\top}\hat{\boldsymbol{\beta}}_g
    -
    \hat{\alpha}_g
  }{\hat{\sigma}_g}.
\]
Within this standardized space, platform effects are estimated separately for AHBA
and GTEx.  The raw ComBat estimates are the batch-specific mean and variance of
$s_{ijg}$,
\[
  \hat{\gamma}_{ig}
  =
  \frac{1}{n_i}\sum_{j=1}^{n_i} s_{ijg},
  \qquad
  \hat{\delta}_{ig}
  =
  \operatorname{Var}_{j=1,\ldots,n_i}(s_{ijg}),
\]
where $n_i$ is the number of rows in batch $i$.  Following the empirical-Bayes
motivation of ComBat, these gene-wise batch estimates are stabilized by shrinkage
across genes, yielding posterior batch adjustments $\hat{\gamma}^{*}_{ig}$ and
$\hat{\delta}^{*}_{ig}$.  In our implementation, this shrinkage is a sample-size
weighted average between the gene-specific batch estimate and the corresponding
batch-wide average across genes, with variance shrinkage performed on the
log-variance scale.  The standardized expression is then corrected by removing the
estimated additive shift and multiplicative scale,
\[
  s^{adj}_{ijg}
  =
  \frac{s_{ijg}-\hat{\gamma}^{*}_{ig}}
       {\sqrt{\hat{\delta}^{*}_{ig}}}.
\]
The corrected expression is reconstructed on the original gene scale while restoring
the fitted covariate contribution:
\[
  y^{corr}_{ijg}
  =
  s^{adj}_{ijg}\hat{\sigma}_g
  +
  \hat{\alpha}_g
  +
  \mathbf{x}_{ij}^{\top}\hat{\boldsymbol{\beta}}_g .
\]

Because AHBA and GTEx do not sample the same biological object, the implementation
adds a final overlap-based calibration after ComBat correction.  For each gene,
corrected AHBA and GTEx values are aggregated to parcels observed in both datasets,
and a gene-wise affine map is fit from corrected GTEx parcel means to corrected AHBA
parcel means:
\[
  \bar{y}^{corr}_{A rg}
  \approx
  a_g \bar{y}^{corr}_{G rg} + d_g ,
\]
where $r$ indexes shared parcels.  The final harmonized values are therefore
\[
  y^{(h)}_{ijg}
  =
  \begin{cases}
    y^{corr}_{ijg}, & i=A,\\
    a_g y^{corr}_{ijg} + d_g, & i=G.
  \end{cases}
\]
Thus AHBA defines the reference harmonized atlas scale, while GTEx observations are
first ComBat-corrected and then calibrated onto that reference scale.  The resulting
$X_A^{(h)}$ and $X_{s,h}^{obs}$ are the shared inputs used by the naive atlas-fill
baseline, DLAM, and PLAM.

\subsection{Naive Atlas-Fill Baseline}
\label{sec:naive-fill}
The same harmonized domain also supports a simple atlas-prior baseline.  For a GTEx
subject $s$, the naive fill predictor preserves all observed harmonized parcels and
fills every unobserved parcel directly from the harmonized AHBA atlas:
\begin{equation}
  \label{eq:naive-fill-unobs}
  \hat{x}_{srg}^{\text{naive}} = x_{A,rg}^{(h)}, \qquad r \notin \Omega_s,
\end{equation}
\begin{equation}
  \label{eq:naive-fill-obs}
  \hat{x}_{srg}^{\text{naive}} = x_{srg}^{(G,h)}, \qquad r \in \Omega_s.
\end{equation}
This baseline is not an alignment model: it contains no subject-specific latent
structure beyond the exact preservation of observed parcels.  It serves as a
harmonized atlas-prior reference for both quantitative benchmarking and residual
visualization.

\section{Deterministic Latent Alignment Model (DLAM)}
\subsection{Intuition}
DLAM treats the observed parcels of a subject as
sufficient to estimate a subject-specific latent expression basis that is spatially
structured.  That basis is aligned to an atlas reference basis by an affine change of
coordinates.  After the subject-specific latent spatial trajectory has been transported
into the atlas frame, it is interpolated over unobserved parcels and decoded back into
gene space.  The procedure is modular: harmonization, latent decomposition, affine
alignment, interpolation, ridge regression, and decoding are estimated sequentially as
point estimates.

\subsection{Joint PLS decomposition}
\label{sec:det-pls}

The key idea of DLAM is to factorize the subject's harmonized
expression matrix $X_{s,h}^{obs}$ into two components: one that is specifically
aligned with the spatial distribution of the observed parcels, and one that captures
the residual gene-expression variation.  This is achieved by fitting a joint partial
least squares (PLS) decomposition that simultaneously decomposes expression and
spatial covariates, finding latent directions of maximum covariance between the two.

Formally, let $\tilde X = X_{s,h}^{obs}$ and $\tilde Y = Y_{\Omega_s}$ denote the
centered expression and spatial matrices at the observed parcels.  PLS finds successive
pairs of weight vectors $(p_k,\, c_k)$ with $p_k\in\mathbb{R}^G$ and
$c_k\in\mathbb{R}^3$ solving
\begin{equation}
  \label{eq:pls-objective}
  (p_k, c_k) = \arg\max_{\|p\|=\|c\|=1} \operatorname{Cov}\!\left(\tilde X\, p,\; \tilde Y\, c\right),
\end{equation}
deflating the matrices between components.  The resulting $K$ component pairs yield
latent score matrices
\begin{equation}
  \label{eq:det-expression-decomp}
  X_{s,h}^{obs} \approx Z_s W_s^\top,
  \qquad
  Z_s \in \mathbb{R}^{|\Omega_s|\times K},\;
  W_s \in \mathbb{R}^{G\times K},
\end{equation}
\begin{equation}
  \label{eq:det-spatial-decomp}
  Y_{\Omega_s} \approx S_s C_s^\top,
  \qquad
  S_s \in \mathbb{R}^{|\Omega_s|\times K},\;
  C_s \in \mathbb{R}^{3\times K},
\end{equation}
where $Z_s$ are latent expression scores and $S_s$ are latent spatial scores.  By
construction of the PLS criterion, each column $z_{s,\cdot,k}$ of $Z_s$ is the linear
combination of genes that maximally co-varies with the corresponding spatial score
$s_{s,\cdot,k}$.  This means $Z_s$ captures precisely the spatially-structured
component of expression: it is the part of $X_{s,h}^{obs}$ that can be predicted
from spatial location.

The atlas reference decomposition is obtained analogously on $X_A^{(h)}$:
\begin{equation}
  \label{eq:det-atlas-decomp}
  X_A^{(h)} \approx Z_A^{det}\!\left(W_A^{det}\right)^\top,
  \qquad
  Z_A^{det}\in\mathbb{R}^{R\times K},\;
  W_A^{det}\in\mathbb{R}^{G\times K}.
\end{equation}

\subsection{Affine alignment, interpolation, and decoding}
\label{sec:det-align}

Because the subject and atlas PLS decompositions are fit independently, the resulting
latent spaces are in different coordinate systems.  The subject latent expression
scores $Z_s$ are aligned to the atlas reference scores $Z_A^{det}$ on the observed
parcels by solving a regularized affine regression:
\begin{equation}
  \label{eq:det-affine-fit}
  (B_s,b_s)
  =
  \arg\min_{B,b}
  \left\|
  Z_s B + \mathbf{1}b^\top - Z_{A,\Omega_s}^{det}
  \right\|_F^2
  +
  \lambda_B \|B\|_F^2.
\end{equation}
The fitted affine map is then applied to both the latent expression and spatial
scores:
\begin{equation}
  \label{eq:det-affine-map}
  Z'_s = Z_s B_s + \mathbf{1}b_s^\top,
\end{equation}
\begin{equation}
  \label{eq:det-spatial-transport}
  S'_s = S_s B_s.
\end{equation}

\paragraph{Why the same $B_s$ applies to both $Z_s$ and $S_s$.}
The PLS decomposition in Eqs.~\eqref{eq:det-expression-decomp}--\eqref{eq:det-spatial-decomp}
operates in a \emph{shared} latent space: both $Z_s$ and $S_s$ are projections of
their respective data matrices onto the same set of $K$ latent directions, chosen to
maximize cross-covariance.  Concretely, one can write
$z_{s,r} = X_{s,r}^{(h)} W_s$ and $s_{s,r} = y_{\Omega_s,r} C_s$, where the loading
directions $W_s$ and $C_s$ are jointly determined by the PLS criterion
\eqref{eq:pls-objective}.  The alignment matrix $B_s$ represents a linear change of
basis in this shared PLS latent space.  Applying it consistently to $Z_s$ and $S_s$
is therefore necessary to maintain the correspondence between expression modes and
spatial modes that was established by PLS: if the expression basis is rotated without
rotating the spatial basis, the relationship between spatial position and predicted
expression is lost.

The aligned latent spatial coordinates are extended to all parcels by an RBF
interpolator
\begin{equation}
  \label{eq:det-rbf}
  \hat{S}'_s(r) = f_s(y_r), \qquad r=1,\dots,R,
\end{equation}
after which a ridge regression maps interpolated latent spatial coordinates back to
aligned latent expression coordinates,
\begin{equation}
  \label{eq:det-ridge-regression}
  \hat{Z}'_s = \hat{S}'_s M_s + \mathbf{1}\beta_s^\top.
\end{equation}
The aligned coordinates are then returned to the subject-specific score system,
\begin{equation}
  \label{eq:det-inverse-affine}
  \hat{Z}_s = \left(\hat{Z}'_s - \mathbf{1}b_s^\top\right) B_s^{-1},
\end{equation}
and decoded into harmonized expression via the subject-specific gene loadings,
\begin{equation}
  \label{eq:det-decode}
  \hat{X}_{s,h} = \hat{Z}_s W_s^\top.
\end{equation}
Observed parcels are preserved exactly by overwrite,
\begin{equation}
  \label{eq:det-overwrite}
  \hat{x}_{srg}^{(h)} = x_{srg}^{(G,h)}, \qquad r\in\Omega_s.
\end{equation}
The deterministic implementation additionally restricts anatomically implausible
borrowing during subject-to-parcel assignment by preferring anchors in the same
hemisphere and brain area before falling back to purely distance-based matches.

\subsection{Initialization}
Initialization is direct because each stage is either closed-form or a standard
regression problem.  The shared harmonizer in Eq.~\eqref{eq:shared-harmonization} is
fit first.  The atlas reference decomposition in Eq.~\eqref{eq:det-atlas-decomp} and
the subject-specific PLS decomposition in
Eqs.~\eqref{eq:det-expression-decomp}--\eqref{eq:det-spatial-decomp} are then fit
from the harmonized matrices.  The affine parameters $(B_s,b_s)$ are initialized by
least squares on the overlap objective in Eq.~\eqref{eq:det-affine-fit}.  The RBF
interpolator in Eq.~\eqref{eq:det-rbf} and the ridge regression in
Eq.~\eqref{eq:det-ridge-regression} are fit on the observed parcels only.  No latent
state is propagated across subjects.

\begin{algorithm}[H]
\caption{Deterministic alignment for subject $s$}
\begin{algorithmic}[1]
\Require Atlas expression $X_A$, subject observations $X_s^{(G)}$ on $\Omega_s$,
         parcel centroids $Y$, latent dimension $K$
\Ensure Completed harmonized atlas $\hat{X}_{s,h}$
\State Fit the shared harmonizer and compute $X_A^{(h)}$ and $X_{s,h}^{obs}$
       using Eq.~\eqref{eq:shared-harmonization}
\State Fit the atlas reference latent decomposition in Eq.~\eqref{eq:det-atlas-decomp}
\State Fit the subject-specific joint PLS decompositions in
       Eqs.~\eqref{eq:det-expression-decomp} and~\eqref{eq:det-spatial-decomp}
\State Estimate $(B_s,b_s)$ by the affine alignment objective in
       Eq.~\eqref{eq:det-affine-fit}
\State Compute aligned subject coordinates by Eqs.~\eqref{eq:det-affine-map}
       and~\eqref{eq:det-spatial-transport}
\State Fit the RBF interpolator and evaluate $\hat{S}'_s(r)$ for all parcels
       using Eq.~\eqref{eq:det-rbf}
\State Fit ridge regression from observed $S'_s$ to observed $Z'_s$ and predict
       $\hat{Z}'_s$ using Eq.~\eqref{eq:det-ridge-regression}
\State Map predicted aligned coordinates back to the subject-specific system with
       Eq.~\eqref{eq:det-inverse-affine}
\State Decode to harmonized expression using Eq.~\eqref{eq:det-decode}
\State Overwrite observed parcels using Eq.~\eqref{eq:det-overwrite}
\State \Return $\hat{X}_{s,h}$
\end{algorithmic}
\end{algorithm}

\section{Probabilistic Latent Alignment Model (PLAM)}
\subsection{Intuition}
PLAM takes the opposite view: instead of estimating a
subject-specific basis first and then transporting it to the atlas, it treats the atlas
as the primary latent geometry and represents each subject as a smooth, aligned
deviation from that geometry.  Subject-specific variation is captured by three
components: an orthogonal alignment matrix, a spatially smooth latent deviation field,
and subject-specific gene-wise calibration parameters.  Unobserved parcels are inferred
by conditioning the deviation field on the observed parcels and then decoding the
completed latent state into harmonized expression.

\subsection{Atlas factor model and subject-specific latent state}
The harmonized atlas is represented by a low-rank factor model,
\begin{equation}
  \label{eq:prob-atlas-factor}
  X_A^{(h)} \approx \mathbf{1}\alpha^\top + Z_A W^\top,
\end{equation}
where $Z_A\in\mathbb{R}^{R\times K}$ contains atlas latent expression coordinates,
$W\in\mathbb{R}^{G\times K}$ is a gene loading matrix, and $\alpha\in\mathbb{R}^G$ is
a gene-wise intercept vector.  This is the generative atlas model: each parcel's
expression profile is explained by a $K$-dimensional latent code $z_{A,r}$ projected
through the shared decoder $W$, plus a gene mean $\alpha$.  It is equivalent to
probabilistic PCA with a fixed mean \cite{tipping1999ppca}.

Let $z_{A,r}$ denote row $r$ of $Z_A$ and $w_g$ denote row $g$ of $W$.  For subject
$s$ and parcel $r$, the latent state is
\begin{equation}
  \label{eq:prob-latent-state}
  z_{s,r} = \left(z_{A,r} + \delta_{s,r}\right)Q_s,
\end{equation}
where $\delta_{s,r}\in\mathbb{R}^K$ is a subject-specific latent deviation from the
atlas at parcel $r$, and $Q_s\in O(K)$ is an orthogonal alignment matrix.  This
decomposition cleanly separates three sources of variation: (i)~the atlas coordinate
$z_{A,r}$, shared across all subjects; (ii)~the subject-specific deviation
$\delta_{s,r}$, which is spatially smooth but subject-dependent; and (iii)~the global
rotation $Q_s$, which captures a linear reorientation of the subject's expression axes
relative to the atlas.  The corresponding platform-free latent expression is
\begin{equation}
  \label{eq:prob-latent-expression}
  x_{srg}^{(*)} = \alpha_g + z_{s,r}^\top w_g.
\end{equation}

\subsection{Subject-specific calibration}
\label{sec:calibration}

After the shared preprocessing harmonization in Eq.~\eqref{eq:shared-harmonization},
residual platform mismatch may persist at the subject level.  PLAM
allows additional gene-wise subject-specific calibration in the observation model:
\begin{equation}
  \label{eq:prob-observation-model}
  x_{srg}^{(G,h)} = a_{s,g}\,x_{srg}^{(*)} + b_{s,g} + \varepsilon_{srg},
  \qquad r\in\Omega_s.
\end{equation}
This calibration layer is structurally equivalent to the location-scale correction in
ComBat \cite{johnson2007combat}, applied at the individual subject level within the
latent observation model rather than at the cohort level in gene space.  While the
global harmonizer in Section~\ref{sec:harmonization-combat} corrects for systematic
platform differences shared across all subjects, the coefficients $(a_{s,g}, b_{s,g})$
absorb residual subject-specific offsets that remain after global correction.

Given the current latent states, the calibration parameters are updated by regularized
regression,
\begin{equation}
  \label{eq:prob-calibration-update}
  \min_{a_{s,g},\,b_{s,g}}
  \sum_{r\in\Omega_s}
  \left(
  x_{srg}^{(G,h)} - a_{s,g}\tilde{x}_{srg}^{(*)} - b_{s,g}
  \right)^2
  +
  \lambda_a (a_{s,g}-a_{0,g})^2
  +
  \lambda_b (b_{s,g}-b_{0,g})^2,
\end{equation}
where $\tilde{x}_{srg}^{(*)}=\alpha_g + z_{s,r}^\top w_g$ and $(a_{0,g},b_{0,g})$
are global calibration targets.  In the neutral initialization, $a_{0,g}=1$ and
$b_{0,g}=0$, corresponding to the prior belief that the global harmonizer has already
removed most platform effects.

\subsection{Full joint distribution}
\label{sec:joint}

The complete PLAM model is summarized by the following factored joint
distribution:
\begin{equation}
  \label{eq:prob-joint}
  \begin{aligned}
  &p\!\left(
  X_A^{(h)}, \{X_s^{(G,h)}\}_{s=1}^S,
  Z_A, W, \alpha,
  \{\Delta_s,Q_s,a_s,b_s\}_{s=1}^S
  \right) \\
  &\quad =
  p(X_A^{(h)}\mid Z_A,W,\alpha)\,
  p(Z_A)\,
  p(W)\,
  p(\alpha) \\
  &\qquad\times
  \prod_{s=1}^S
  p(Q_s)\,
  p(\Delta_s\mid Y)\,
  p(a_s,b_s)\,
  p(X_s^{(G,h)}\mid Z_A,\Delta_s,Q_s,W,\alpha,a_s,b_s).
  \end{aligned}
\end{equation}
The explicit distributional assumptions for each factor are as follows.

\paragraph{Atlas likelihood.}
\[
  p(X_A^{(h)}\mid Z_A,W,\alpha)
  = \prod_{r=1}^R\prod_{g=1}^G
    \mathcal{N}\!\left(x_{A,rg}^{(h)};\;\alpha_g + z_{A,r}^\top w_g,\;\sigma_A^2\right).
\]
This is a Gaussian factor model for the atlas, identical to probabilistic PCA
\cite{tipping1999ppca} with a fixed intercept $\alpha$.

\paragraph{Atlas latent prior.}
\[
  p(Z_A) = \prod_{r=1}^R \mathcal{N}(z_{A,r};\,\mathbf{0},\,I_K).
\]
A standard spherical Gaussian prior on atlas latent codes.

\paragraph{Decoder and intercept priors.}
\[
  p(W) = \prod_{g=1}^G \mathcal{N}(w_g;\,\mathbf{0},\,\tau_W^2 I_K),
  \qquad
  p(\alpha) = \prod_{g=1}^G \mathcal{N}(\alpha_g;\,0,\,\tau_\alpha^2).
\]
Gaussian shrinkage on gene loadings and gene means.

\paragraph{Orthogonal alignment.}
\[
  p(Q_s) = \mathrm{Uniform}\!\left(O(K)\right).
\]
$Q_s$ is constrained to the group of $K\times K$ orthogonal matrices and estimated by
solving a Procrustes problem (Section~\ref{sec:procrustes}).

\paragraph{Spatial GP prior on deviations.}
\[
  p(\Delta_s \mid Y) = \prod_{k=1}^K\,
  \mathcal{GP}\!\left(\delta_{s,\cdot,k};\;0,\;\kappa_k(\cdot,\cdot)\right),
\]
where $\kappa_k$ is a squared-exponential kernel defined over parcel centroids
(Section~\ref{sec:gp-completion}).  This prior encodes the assumption that
subject-specific deviations from the atlas are spatially smooth.

\paragraph{Calibration prior.}
\[
  p(a_s,b_s)
  = \prod_{g=1}^G
    \mathcal{N}(a_{s,g};\,a_{0,g},\,\lambda_a^{-1})\,
    \mathcal{N}(b_{s,g};\,b_{0,g},\,\lambda_b^{-1}).
\]
Gaussian regularization toward the identity transformation $(a_{0,g}=1,\,b_{0,g}=0)$.

\paragraph{Subject likelihood.}
\[
  p(X_s^{(G,h)}\mid \cdots)
  = \prod_{r\in\Omega_s}\prod_{g=1}^G
    \mathcal{N}\!\left(x_{srg}^{(G,h)};\;
    a_{s,g}\bigl(\alpha_g + z_{s,r}^\top w_g\bigr) + b_{s,g},\;
    \sigma_{s,g}^2\right),
\]
where $z_{s,r} = (z_{A,r}+\delta_{s,r})Q_s$ from Eq.~\eqref{eq:prob-latent-state}.

\subsection{From unweighted to weighted latent regression}
The latent-state update is easiest to motivate from ordinary least squares.  If every
gene were treated as equally reliable and all residuals were penalized equally, the
subject-specific latent coordinate at parcel $r$ would solve
\begin{equation}
  \label{eq:prob-unweighted-latent-update}
  \min_{z}
  \sum_{g=1}^{G}
  \left(
  x_{srg}^{(G,h)} -
  a_{s,g}\left(\alpha_g + z^\top w_g\right) -
  b_{s,g}
  \right)^2
  +
  \lambda_z \left\| z - z_{A,r}Q_s \right\|_2^2.
\end{equation}
This objective balances fitting the observed gene expression against staying close to
the atlas-aligned coordinate $z_{A,r}Q_s$.

In practice, some residuals are unusually large and some genes are systematically
noisier than others.  Let
\[
  e_{srg}
  = x_{srg}^{(G,h)}
  - a_{s,g}\left(\alpha_g + z_{s,r}^\top w_g\right)
  - b_{s,g}
\]
denote the current residual.  A Student-$t$-inspired robust weight is
\begin{equation}
  \label{eq:prob-robust-weight}
  \omega_{srg}^{robust}
  =
  \frac{\nu + 1}{\nu + e_{srg}^2 / \hat{\sigma}_{s,g}^2},
\end{equation}
and a gene-wise heteroscedastic weight is
\begin{equation}
  \label{eq:prob-hetero-weight}
  \omega_{sg}^{hetero}
  =
  \left(\hat{\sigma}_{s,g}^2 + \epsilon_\omega\right)^{-1},
  \qquad
  \hat{\sigma}_{s,g}^2
  =
  \frac{1}{|\Omega_s|}
  \sum_{r\in\Omega_s}
  e_{srg}^2.
\end{equation}
The combined weight is
\begin{equation}
  \label{eq:prob-total-weight}
  \omega_{srg} = \omega_{srg}^{robust}\,\omega_{sg}^{hetero}.
\end{equation}
The implemented latent-state update is then the weighted ridge regression
\begin{equation}
  \label{eq:prob-latent-update}
  \min_{z}
  \sum_{g=1}^{G}
  \omega_{srg}
  \left(
  x_{srg}^{(G,h)} -
  a_{s,g}\left(\alpha_g + z^\top w_g\right) -
  b_{s,g}
  \right)^2
  +
  \lambda_z \left\| z - z_{A,r}Q_s \right\|_2^2.
\end{equation}

\paragraph{Connection to M-estimation and IRLS.}
The combined weight $\omega_{srg}$ corresponds to M-estimation
\cite{huber1981robust}: $\omega_{srg}^{robust}$ is proportional to the influence
function of a Student-$t$ distribution, downweighting outlier genes at individual
observations, while $\omega_{sg}^{hetero}$ accounts for gene-level heteroscedasticity.
Fixing $\omega_{srg}$ and solving Eq.~\eqref{eq:prob-latent-update}, then updating
$\omega_{srg}$ given the new $z_{s,r}$, is precisely the iteratively reweighted least
squares (IRLS) algorithm \cite{huber1981robust}.  Each cycle is guaranteed not to
increase the M-estimation objective.

\subsection{Orthogonal Procrustes alignment}
\label{sec:procrustes}

Given the updated observed latent coordinates, the orthogonal alignment matrix $Q_s$
is refreshed by solving the classical orthogonal Procrustes problem
\cite{schonemann1966procrustes}:
\begin{equation}
  \label{eq:prob-procrustes}
  Q_s^\star
  =
  \arg\min_{Q\in O(K)}
  \left\|
  Z_{A,\Omega_s}Q - Z_{s,\Omega_s}
  \right\|_F^2.
\end{equation}
This has the closed-form solution $Q_s^\star = UV^\top$, where
$Z_{A,\Omega_s}^\top Z_{s,\Omega_s} = U\Sigma V^\top$ is the singular value
decomposition.

\subsection{Gaussian-process completion of the latent deviation field}
\label{sec:gp-completion}

The Gaussian process is used to interpolate the subject-specific latent deviation
field rather than the expression directly \cite{rasmussen2006gpml}.  After the observed
latent coordinates have been updated, the deviation at observed parcels is obtained by
undoing the orthogonal alignment and subtracting the atlas coordinate:
\begin{equation}
  \label{eq:prob-observed-deviation}
  \delta_{s,r}^{obs} = z_{s,r}Q_s^\top - z_{A,r},
  \qquad r\in\Omega_s.
\end{equation}
For latent dimension $k$, the scalar deviations
$\{\delta_{s,r,k}^{obs}\}_{r\in\Omega_s}$ are modeled as noise-corrupted observations
from a zero-mean Gaussian process with squared-exponential kernel
\begin{equation}
  \label{eq:prob-kernel}
  \kappa_k(y_r,y_{r'})
  =
  \sigma_k^2
  \exp\!\left(
  -\frac{\|y_r-y_{r'}\|_2^2}{2\ell_k^2}
  \right),
\end{equation}
where $\sigma_k^2$ is the marginal signal variance and $\ell_k$ is the spatial
length-scale for dimension $k$.  Let $K_{\Omega,\Omega}^{(k)}$ denote the
$|\Omega_s|\times|\Omega_s|$ kernel matrix evaluated at observed parcel centroids,
$K_{r,\Omega}^{(k)}$ the $1\times|\Omega_s|$ cross-kernel between parcel $r$ and
the observed parcels, and $K_{r,r}^{(k)}$ the prior variance at $r$.  The GP
posterior mean and variance at any parcel $r$ are then
\begin{equation}
  \label{eq:prob-gp-mean}
  \hat{\delta}_{s,r,k}
  =
  K_{r,\Omega}^{(k)}
  \left(
  K_{\Omega,\Omega}^{(k)} + \tau_k^2 I
  \right)^{-1}
  \delta_{s,\Omega,k}^{obs},
\end{equation}
\begin{equation}
  \label{eq:prob-gp-var}
  \mathrm{Var}\!\left(
  \delta_{s,r,k}\mid \Omega_s
  \right)
  =
  K_{r,r}^{(k)}
  -
  K_{r,\Omega}^{(k)}
  \left(
  K_{\Omega,\Omega}^{(k)} + \tau_k^2 I
  \right)^{-1}
  K_{\Omega,r}^{(k)},
\end{equation}
where $\tau_k^2$ is an observation noise variance \cite{rasmussen2006gpml}.  The
posterior mean $\hat\delta_{s,r,k}$ provides the completed deviation at unobserved
parcel $r$, while the posterior variance quantifies how uncertain that completion is
as a function of distance from the observed parcels.

Averaging the posterior variance over latent dimensions gives the parcel-wise average
latent posterior variance,
\begin{equation}
  \label{eq:prob-avg-latent-var}
  \bar{\sigma}_{s,r}^{2}
  =
  \frac{1}{K}
  \sum_{k=1}^{K}
  \mathrm{Var}\!\left(
  \delta_{s,r,k}\mid \Omega_s
  \right).
\end{equation}
The completed latent coordinate is reconstructed as
\begin{equation}
  \label{eq:prob-completed-latent}
  \hat{z}_{s,r}
  =
  \left(
  z_{A,r} + \hat{\delta}_{s,r}
  \right)Q_s.
\end{equation}
The completed harmonized expression is decoded as
\begin{equation}
  \label{eq:prob-final-decode}
  \hat{x}_{srg}^{(h)}
  =
  a_{s,g}\!\left(
  \alpha_g + \hat{z}_{s,r}^\top w_g
  \right)
  +
  b_{s,g},
\end{equation}
and observed parcels are preserved exactly by
\begin{equation}
  \label{eq:prob-overwrite}
  \hat{x}_{srg}^{(h)} = x_{srg}^{(G,h)}, \qquad r\in\Omega_s.
\end{equation}

\subsection{Optional uncertainty shrinkage}
The broader probabilistic family also permits the completed subject atlas to be
shrunk toward an atlas prior according to spatial distance and average latent posterior
variance.  If $\tilde{d}_{s,r}$ is a scaled distance to the nearest observed parcel
and $\widetilde{\bar{\sigma}}_{s,r}^{2}$ is the scaled average latent posterior
variance from Eq.~\eqref{eq:prob-avg-latent-var}, then one possible shrinkage score is
\[
  \operatorname{score}_{s,r}
  =
  \gamma_d \tilde{d}_{s,r}
  +
  \gamma_u \widetilde{\bar{\sigma}}_{s,r}^{2},
\]
with logistic weight
\[
  w_{s,r}
  =
  \sigma\!\left(
  \frac{\operatorname{score}_{s,r}-m_0}{\tau}
  \right),
  \qquad
  \hat{x}_{s,r}^{final}
  =
  (1-w_{s,r})\hat{x}_{s,r}^{model}
  +
  w_{s,r}\hat{x}_{r}^{prior}.
\]
This extension belongs to the general model family but is inactive in the particular
probabilistic analysis presented here.

\subsection{Initialization}
Initialization is explicit because PLAM is fit by alternating
updates.  After the shared preprocessing harmonization in
Eq.~\eqref{eq:shared-harmonization}, the atlas latent coordinates $Z_A$ are
initialized by a rank-$K$ singular value decomposition of centered $X_A^{(h)}$, and
the decoder parameters $(W,\alpha)$ are fit by ridge regression under
Eq.~\eqref{eq:prob-atlas-factor}.  For each subject, the orthogonal alignment is
initialized as $Q_s=I$, the observed latent coordinates are initialized as the atlas
coordinates at the observed parcels, $z_{s,r}=z_{A,r}$ for $r\in\Omega_s$, the
calibration parameters are initialized as $a_{s,g}=1$ and $b_{s,g}=0$, and the
combined weights are initialized as $\omega_{srg}=1$.  These choices correspond to
the neutral hypothesis that the subject initially coincides with the atlas in the
harmonized domain.

\begin{algorithm}[H]
\caption{Unified probabilistic alignment for subject $s$}
\begin{algorithmic}[1]
\Require Atlas expression $X_A$, subject observations $X_s^{(G)}$ on $\Omega_s$,
         parcel centroids $Y$, latent dimension $K$, number of iterations $T$
\Ensure Completed harmonized atlas $\hat{X}_{s,h}$ and parcel-wise average latent
        posterior variance $\bar{\sigma}_{s,r}^2$
\State Fit the shared harmonizer and compute $X_A^{(h)}$ and $X_{s,h}^{obs}$
       using Eq.~\eqref{eq:shared-harmonization}
\State Initialize $Z_A$, $W$, and $\alpha$ under the atlas factor model in
       Eq.~\eqref{eq:prob-atlas-factor}
\State Initialize $Q_s \gets I$, $z_{s,r} \gets z_{A,r}$ for $r\in\Omega_s$,
       $a_{s,g}\gets 1$, $b_{s,g}\gets 0$, and $\omega_{srg}\gets 1$
\For{$t=1,\dots,T$}
    \For{each gene $g$}
        \State Update $(a_{s,g},b_{s,g})$ by solving the regularized regression in
               Eq.~\eqref{eq:prob-calibration-update}
    \EndFor
    \For{each observed parcel $r\in\Omega_s$}
        \State Form residuals $e_{srg}$ and update $\omega_{srg}^{robust}$,
               $\omega_{sg}^{hetero}$, and $\omega_{srg}$ using
               Eqs.~\eqref{eq:prob-robust-weight}--\eqref{eq:prob-total-weight}
        \State Update $z_{s,r}$ by solving the weighted latent regression in
               Eq.~\eqref{eq:prob-latent-update}
    \EndFor
    \State Update $Q_s$ by the orthogonal Procrustes step in
           Eq.~\eqref{eq:prob-procrustes}
\EndFor
\State Form observed deviations $\delta_{s,r}^{obs}$ using
       Eq.~\eqref{eq:prob-observed-deviation}
\For{each latent dimension $k=1,\dots,K$}
    \State Compute the Gaussian-process posterior mean and variance using
           Eqs.~\eqref{eq:prob-gp-mean} and~\eqref{eq:prob-gp-var}
\EndFor
\State Compute $\bar{\sigma}_{s,r}^{2}$ using Eq.~\eqref{eq:prob-avg-latent-var}
       and the completed latent coordinates $\hat{z}_{s,r}$ using
       Eq.~\eqref{eq:prob-completed-latent}
\State Decode to harmonized expression using Eq.~\eqref{eq:prob-final-decode}
\State Overwrite observed parcels using Eq.~\eqref{eq:prob-overwrite}
\State \Return $\hat{X}_{s,h}$ and $\bar{\sigma}_{s,r}^{2}$
\end{algorithmic}
\end{algorithm}

\section{Atlas-Level Diagnostics}
Atlas-level diagnostics are useful because they separate two questions: whether the
harmonization and alignment steps place AHBA and GTEx in a comparable coordinate
system, and whether the imputation step preserves parcel-level structure when sparse
subject data are aggregated into a population atlas.
Figures~\ref{fig:atlas-det-heatmap} and~\ref{fig:atlas-uni-heatmap} show the
progression from raw AHBA and sparse GTEx observations to harmonized and completed
parcel-by-gene atlases.  In both cases, the gap between the sparse observed GTEx
pattern and the completed harmonized atlas is explicitly visible.
Figures~\ref{fig:atlas-naive-heatmap} and~\ref{fig:atlas-naive-scatter} provide the
same diagnostics for the naive atlas-fill baseline, making it possible to compare
learned alignment against direct atlas borrowing in the same harmonized domain.

\begin{figure}[H]
\centering
\includegraphics[width=0.94\textwidth]{figs/atlas_deterministic_heatmap_demo.pdf}
\caption{Atlas-level heatmap diagnostic for DLAM.  Panels show (A) AHBA raw,
(B) sparse GTEx raw, (C) sparse GTEx after harmonization, (D) AHBA harmonized,
(E) the DLAM-completed harmonized GTEx atlas, and (F) the residual difference
relative to the harmonized AHBA atlas.}
\label{fig:atlas-det-heatmap}
\end{figure}

\begin{figure}[H]
\centering
\includegraphics[width=0.94\textwidth]{figs/atlas_unified_heatmap_demo.pdf}
\caption{Atlas-level heatmap diagnostic for PLAM.  Panels show (A) AHBA raw,
(B) sparse GTEx raw, (C) sparse GTEx after harmonization, (D) AHBA harmonized,
(E) the PLAM-completed harmonized GTEx atlas, and (F) the residual difference
relative to the harmonized AHBA atlas.}
\label{fig:atlas-uni-heatmap}
\end{figure}

\begin{figure}[H]
\centering
\includegraphics[width=0.94\textwidth]{figs/atlas_naive_heatmap_demo.pdf}
\caption{Atlas-level heatmap diagnostic for the naive atlas-fill baseline.
Panels show (A) AHBA raw, (B) sparse GTEx raw, (C) sparse GTEx after harmonization,
(D) AHBA harmonized, (E) the naive-fill harmonized GTEx atlas, and (F) the residual
difference relative to the harmonized AHBA atlas.}
\label{fig:atlas-naive-heatmap}
\end{figure}

\begin{figure}[H]
\centering
\begin{subfigure}[t]{0.49\textwidth}
    \centering
    \includegraphics[width=\textwidth]{figs/atlas_deterministic_scatter_region_harmonized.pdf}
\end{subfigure}
\hfill
\begin{subfigure}[t]{0.49\textwidth}
    \centering
    \includegraphics[width=\textwidth]{figs/atlas_deterministic_scatter_gene_harmonized.pdf}
\end{subfigure}
\medskip
\begin{subfigure}[t]{0.49\textwidth}
    \centering
    \includegraphics[width=\textwidth]{figs/atlas_deterministic_scatter_region_corr.pdf}
\end{subfigure}
\hfill
\begin{subfigure}[t]{0.49\textwidth}
    \centering
    \includegraphics[width=\textwidth]{figs/atlas_deterministic_scatter_gene_corr.pdf}
\end{subfigure}
\caption{Atlas-level scatter diagnostics for DLAM.  Panels show (A) parcel identity,
(B) gene identity, (C) parcel-level correlation, and (D) gene-level correlation in the
harmonized atlas space.  Only the most informative parcel and gene subsets are
individually color-coded, with the remaining points shown in gray for context.}
\end{figure}

\begin{figure}[H]
\centering
\begin{subfigure}[t]{0.49\textwidth}
    \centering
    \includegraphics[width=\textwidth]{figs/atlas_unified_scatter_region_harmonized.pdf}
\end{subfigure}
\hfill
\begin{subfigure}[t]{0.49\textwidth}
    \centering
    \includegraphics[width=\textwidth]{figs/atlas_unified_scatter_gene_harmonized.pdf}
\end{subfigure}
\medskip
\begin{subfigure}[t]{0.49\textwidth}
    \centering
    \includegraphics[width=\textwidth]{figs/atlas_unified_scatter_region_corr.pdf}
\end{subfigure}
\hfill
\begin{subfigure}[t]{0.49\textwidth}
    \centering
    \includegraphics[width=\textwidth]{figs/atlas_unified_scatter_gene_corr.pdf}
\end{subfigure}
\caption{Atlas-level scatter diagnostics for PLAM.  Panels show (A) parcel identity,
(B) gene identity, (C) parcel-level correlation, and (D) gene-level correlation in the
harmonized atlas space.  The same visual grammar is used as in the DLAM panel so that
the effect of the two alignment strategies can be compared directly.}
\label{fig:atlas-plam-scatter}
\end{figure}

\begin{figure}[H]
\centering
\begin{subfigure}[t]{0.49\textwidth}
    \centering
    \includegraphics[width=\textwidth]{figs/atlas_naive_scatter_region_harmonized.pdf}
\end{subfigure}
\hfill
\begin{subfigure}[t]{0.49\textwidth}
    \centering
    \includegraphics[width=\textwidth]{figs/atlas_naive_scatter_gene_harmonized.pdf}
\end{subfigure}
\medskip
\begin{subfigure}[t]{0.49\textwidth}
    \centering
    \includegraphics[width=\textwidth]{figs/atlas_naive_scatter_region_corr.pdf}
\end{subfigure}
\hfill
\begin{subfigure}[t]{0.49\textwidth}
    \centering
    \includegraphics[width=\textwidth]{figs/atlas_naive_scatter_gene_corr.pdf}
\end{subfigure}
\caption{Atlas-level scatter diagnostics for the naive atlas-fill baseline.
Panels show (A) parcel identity, (B) gene identity, (C) parcel-level correlation, and
(D) gene-level correlation in the harmonized atlas space.}
\label{fig:atlas-naive-scatter}
\end{figure}

\section{Single-Subject Diagnostics}
Single-subject diagnostics illustrate how each model turns sparse observations into a
completed atlas.  The representative subject shown here is \texttt{GTEX-14ASI}, chosen
because its subject-level PLAM performance lies close to the cohort median
among subjects with at least eight observed parcels.  This makes the visual example
representative of typical behavior rather than an extreme case.

\begin{figure}[H]
\centering
\begin{subfigure}[t]{0.32\textwidth}
    \centering
    \includegraphics[width=\textwidth]{figs/subject_deterministic_heatmap_demo.pdf}
    \caption{DLAM}
\end{subfigure}
\hfill
\begin{subfigure}[t]{0.32\textwidth}
    \centering
    \includegraphics[width=\textwidth]{figs/subject_unified_heatmap_demo.pdf}
    \caption{PLAM}
\end{subfigure}
\hfill
\begin{subfigure}[t]{0.32\textwidth}
    \centering
    \includegraphics[width=\textwidth]{figs/subject_naive_heatmap_demo.pdf}
    \caption{Naive fill}
\end{subfigure}
\caption{Single-subject heatmap diagnostics for the representative subject
\texttt{GTEX-14ASI}.  Panels (A)--(C) compare the transition from sparse observations
to the completed harmonized atlas under DLAM, PLAM, and naive fill in the same visual
format.}
\end{figure}

\begin{figure}[H]
\centering
\begin{subfigure}[t]{0.32\textwidth}
    \centering
    \includegraphics[width=\textwidth]{figs/subject_deterministic_scatter_harmonized.pdf}
    \caption{DLAM}
\end{subfigure}
\hfill
\begin{subfigure}[t]{0.32\textwidth}
    \centering
    \includegraphics[width=\textwidth]{figs/subject_unified_scatter_harmonized.pdf}
    \caption{PLAM}
\end{subfigure}
\hfill
\begin{subfigure}[t]{0.32\textwidth}
    \centering
    \includegraphics[width=\textwidth]{figs/subject_naive_scatter_harmonized.pdf}
    \caption{Naive fill}
\end{subfigure}
\caption{Single-subject scatter diagnostics in harmonized space for the representative
subject.  Panels (A)--(C) show parcel identity, gene identity, parcel-level
correlation, and gene-level correlation for the completed subject atlas under DLAM,
PLAM, and naive fill.}
\end{figure}

\section{Validation Context}
Validation is based on leave-one-region-out cross-validation at the subject level.
For each subject and each observed parcel, the parcel is held out completely, all
fitted quantities are recomputed without it, and the held-out gene vector is predicted
from the remaining parcels.  This protocol is strict: the held-out parcel is excluded
from harmonization, latent fitting, spatial completion, and decoding.  The resulting
diagnostics therefore measure genuine out-of-sample regional generalization rather than
reconstruction of inputs seen during fitting.

\paragraph{Naive baseline.}
Both alignment models are benchmarked against a naive atlas-fill baseline.  For each held-out
parcel $r$ of subject $s$, the predicted expression is simply the corresponding row of
the harmonized AHBA atlas:
\[
  \hat{x}_{srg}^{naive} = x_{A,rg}^{(h)},
  \qquad g = 1,\dots,G.
\]
This baseline captures the information available purely from the atlas spatial prior
without any subject-specific adaptation.  Improvement over this baseline therefore
reflects the additional benefit of fitting each model to the individual subject's
observed parcels.

The leave-one-region-out diagnostics reported below are computed over the full gene
set so that the manuscript-level evaluation is consistent with the main
transcriptome-wide completion setting.  These plots are diagnostic rather than
exhaustive: the aim is to expose how error and correlation vary over subjects, parcels,
and genes, not merely to report a single cohort-level score.

\begin{figure}[H]
\centering
\includegraphics[width=0.95\textwidth]{figs/loro_subject_metrics_allgene.pdf}
\caption{Subject-level leave-one-region-out diagnostics over the full gene set.
Subjects are ordered by PLAM Pearson correlation.  The panels show the subject-wise
coverage indicator, subject-wise Pearson correlation, subject-wise RMSE, and cohort-level
summary bars for DLAM, PLAM, and the naive AHBA-fill baseline.}
\end{figure}

\begin{figure}[H]
\centering
\includegraphics[width=0.95\textwidth]{figs/loro_parcel_metrics_allgene.pdf}
\caption{Parcel-level leave-one-region-out diagnostics over the full gene set.  The
heatmaps summarize mean held-out correlation and RMSE by parcel across the canonical
parcel set; rows with no GTEx holdout events are left blank.  The bar plot shows
parcel-specific improvement relative to the naive AHBA-fill baseline for parcels with
available leave-one-region-out evaluations.}
\end{figure}

\begin{figure}[H]
\centering
\includegraphics[width=0.92\textwidth]{figs/loro_method_summary_bars_allgene.pdf}
\caption{Cohort-level summary of leave-one-region-out diagnostics over the full gene
set.  Mean Pearson correlation, mean RMSE, and mean improvement relative to the naive
AHBA-fill baseline are shown for both atlas-completion models and the naive baseline.}
\end{figure}

\section{Discussion}
DLAM and PLAM share a common scientific aim: to infer a subject-complete parcel-by-gene atlas
from sparse GTEx regional measurements by borrowing structure from AHBA.  Their main
difference lies in what is treated as primary.  DLAM treats each
subject as the starting point, estimates a subject-specific latent basis whose spatial
structure is revealed by PLS, and then transports that basis into the atlas frame.
PLAM treats the atlas as the starting point, defines a global latent
geometry there, and then expresses each subject as a smooth, orthogonally aligned
perturbation of that geometry.

This distinction has practical consequences.  DLAM is modular and
interpretable at each step, which makes it attractive for diagnosing where alignment
or extrapolation may fail.  PLAM is more coherent because
calibration, alignment, latent completion, and uncertainty all live in a single
atlas-referenced latent system.  That coherence is what makes Gaussian-process
deviation conditioning natural in the probabilistic formulation: the GP is not
interpolating gene expression directly, but a low-dimensional deviation field whose
spatial smoothness is biologically and statistically motivated \cite{rasmussen2006gpml}.

The two approaches should therefore not be understood simply as different stages of a
software pipeline.  They represent two different answers to the same problem.  DLAM
provides a transparent latent transport view built from subject-specific components,
whereas PLAM provides a global
generative view in which atlas structure, subject alignment, subject deviation, and
calibration are estimated jointly in a common latent space.

\begin{thebibliography}{99}

\bibitem[Hawrylycz et~al.(2012)]{hawrylycz2012anatomically}
M.~J. Hawrylycz et~al.
\newblock An anatomically comprehensive atlas of the adult human brain transcriptome.
\newblock {\em Nature}, 489:391--399, 2012.

\bibitem[GTEx Consortium(2013)]{gtex2013genotype}
The GTEx Consortium.
\newblock The genotype-tissue expression (GTEx) project.
\newblock {\em Nature Genetics}, 45:580--585, 2013.

\bibitem[GTEx Consortium(2017)]{gtex2017genetic}
The GTEx Consortium.
\newblock Genetic effects on gene expression across human tissues.
\newblock {\em Nature}, 550:204--213, 2017.

\bibitem[GTEx Consortium(2020)]{gtex2020atlas}
The GTEx Consortium.
\newblock The GTEx Consortium atlas of genetic regulatory effects across human tissues.
\newblock {\em Science}, 369:1318--1330, 2020.

\bibitem[Arnatkevi{\v{c}}i{\={u}}t{\.{e}} et~al.(2019)]{arnatkeviciute2019practical}
A. Arnatkevi{\v{c}}i{\={u}}t{\.{e}}, B.~D. Fulcher, and A. Fornito.
\newblock A practical guide to linking brain-wide gene expression and neuroimaging data.
\newblock {\em NeuroImage}, 189:353--367, 2019.

\bibitem[Markello et~al.(2021)]{markello2021abagen}
R.~D. Markello et~al.
\newblock Standardizing workflows in imaging transcriptomics with the abagen toolbox.
\newblock {\em eLife}, 10:e72129, 2021.

\bibitem[Richiardi et~al.(2015)]{richiardi2015correlated}
J. Richiardi et~al.
\newblock Correlated gene expression supports synchronous activity in brain networks.
\newblock {\em Science}, 348:1241--1244, 2015.

\bibitem[Fornito et~al.(2019)]{fornito2019bridging}
A. Fornito, A. Arnatkevi{\v{c}}i{\={u}}t{\.{e}}, and B.~D. Fulcher.
\newblock Bridging the gap between connectome and transcriptome.
\newblock {\em Trends in Cognitive Sciences}, 23:34--50, 2019.

\bibitem[Burt et~al.(2018)]{burt2018hierarchy}
J.~B. Burt et~al.
\newblock Hierarchy of transcriptomic specialization across human cortex captured by
structural neuroimaging topography.
\newblock {\em Nature Neuroscience}, 21:1251--1259, 2018.

\bibitem[Vogel et~al.(2023)]{vogel2023molecular}
J.~W. Vogel et~al.
\newblock Molecular gradients and cross-atlas spatial transcriptomics alignment.
\newblock {\em Proceedings of the National Academy of Sciences},
120:e2219137121, 2023.

\bibitem[Geladi and Kowalski(1986)]{geladi1986pls}
P. Geladi and B.~R. Kowalski.
\newblock Partial least-squares regression: a tutorial.
\newblock {\em Analytica Chimica Acta}, 185:1--17, 1986.

\bibitem[de~Jong(1993)]{dejong1993simpls}
S. de~Jong.
\newblock SIMPLS: an alternative approach to partial least squares regression.
\newblock {\em Journal of Chemometrics}, 7:251--263, 1993.

\bibitem[Wold et~al.(2001)]{wold2001pls}
S. Wold, M. Sj\"ostr\"om, and L. Eriksson.
\newblock PLS-regression: a basic tool of chemometrics.
\newblock {\em Chemometrics and Intelligent Laboratory Systems}, 58:109--130, 2001.

\bibitem[Sch\"onemann(1966)]{schonemann1966procrustes}
P.~H. Sch\"onemann.
\newblock A generalized solution of the orthogonal Procrustes problem.
\newblock {\em Psychometrika}, 31:1--10, 1966.

\bibitem[Johnson et~al.(2007)]{johnson2007combat}
W.~E. Johnson, C. Li, and A. Rabinovic.
\newblock Adjusting batch effects in microarray expression data using empirical Bayes
methods.
\newblock {\em Biostatistics}, 8:118--127, 2007.

\bibitem[Bolstad et~al.(2003)]{bolstad2003normalization}
B.~M. Bolstad et~al.
\newblock A comparison of normalization methods for high density oligonucleotide array
data based on variance and bias.
\newblock {\em Bioinformatics}, 19:185--193, 2003.

\bibitem[Tipping and Bishop(1999)]{tipping1999ppca}
M.~E. Tipping and C.~M. Bishop.
\newblock Probabilistic principal component analysis.
\newblock {\em Journal of the Royal Statistical Society, Series B}, 61:611--622, 1999.

\bibitem[Huber(1981)]{huber1981robust}
P.~J. Huber.
\newblock {\em Robust Statistics}.
\newblock Wiley, New York, 1981.

\bibitem[Rasmussen and Williams(2006)]{rasmussen2006gpml}
C.~E. Rasmussen and C.~K.~I. Williams.
\newblock {\em Gaussian Processes for Machine Learning}.
\newblock MIT Press, 2006.

\bibitem[Cressie(1993)]{cressie1993statistics}
N.~A.~C. Cressie.
\newblock {\em Statistics for Spatial Data}.
\newblock Wiley, 1993.

\bibitem[Stuart et~al.(2019)]{stuart2019comprehensive}
T. Stuart et~al.
\newblock Comprehensive integration of single-cell data.
\newblock {\em Cell}, 177:1888--1902, 2019.

\bibitem[Korsunsky et~al.(2019)]{korsunsky2019harmony}
I. Korsunsky et~al.
\newblock Fast, sensitive and accurate integration of single-cell data with Harmony.
\newblock {\em Nature Methods}, 16:1289--1296, 2019.

\bibitem[Lotfollahi et~al.(2022)]{lotfollahi2022scarches}
M. Lotfollahi et~al.
\newblock Mapping single-cell data to reference atlases by transfer learning.
\newblock {\em Nature Biotechnology}, 40:121--130, 2022.

\bibitem[Li et~al.(2022)]{li2022tangram}
B. Li et~al.
\newblock Integrating spatial transcriptomics and single-cell transcriptomics with
Tangram.
\newblock {\em Nature Methods}, 19:317--324, 2022.

\bibitem[Kleshchevnikov et~al.(2022)]{kleshchevnikov2022cell2location}
V. Kleshchevnikov et~al.
\newblock Cell2location maps fine-grained cell types in spatial transcriptomics.
\newblock {\em Nature Biotechnology}, 40:661--671, 2022.

\end{thebibliography}

\end{document}
